% =====================================================================
%  P3.  The exact second moment of sum_{n<N} Lambda(n) mu(N-n),
%       and why Chowla's conjecture does not control it.
%
%  Deployment version.  Written to deploy/SPEC.md.
%  Compile:  pdflatex P3-wall-second-moment.tex   (twice)
% =====================================================================
\documentclass[11pt,a4paper]{article}

\usepackage{amsmath}
\usepackage{amssymb}
\usepackage{amsthm}
\usepackage[margin=1in]{geometry}
\usepackage[colorlinks=true,linkcolor=blue,citecolor=blue]{hyperref}

\theoremstyle{plain}
\newtheorem{theorem}{Theorem}
\newtheorem{corollary}[theorem]{Corollary}
\newtheorem{proposition}[theorem]{Proposition}
\newtheorem{lemma}[theorem]{Lemma}
\newtheorem{conjecture}[theorem]{Conjecture}
\theoremstyle{definition}
\newtheorem{observation}[theorem]{Observation}
\newtheorem{measurement}[theorem]{Measurement}
\theoremstyle{remark}
\newtheorem{note}[theorem]{Note}

\renewcommand{\SS}{\mathfrak{S}}
\newcommand{\AAA}{\mathfrak{A}}

% typographic slack: let TeX stretch a line rather than run into the margin
\emergencystretch=3em

\title{The exact second moment of
$\sum_{n<N}\Lambda(n)\mu(N-n)$,\\
and why Chowla's conjecture does not control it}
\author{}
\date{}

\begin{document}
\maketitle

\begin{abstract}
The Huang--Li reduction of binary Goldbach to Elliott--Halberstam
\cite{HL}, once its demand side is evaluated \cite{P1}, rests on the
single scalar
\[
  C(N)\;=\;\sum_{n<N}\Lambda(n)\mu(N-n)\;=\;o(N).
\]
This paper establishes what is exactly computable about that scalar.
Its second moment has a closed form,
$V(N)=\sum_{v<N}\mu^2(v)\Lambda(N-v)^2\sim\AAA(N)N\log N$, whose local
factor is $\AAA(N)=\prod_{q\nmid N}(1-1/(q(q-1)))$ and \emph{not} the
Hardy--Littlewood singular series (Proposition~\ref{prop:V}); the same
constant is, not coincidentally, the deficiency in Huang--Li's own
lower bound. The aggregate second moment of the truncated convolution
is an exact unconditional identity in terms of the prime-pair sum and
the binary Chowla sum (Lemma~\ref{lem:MP}), the truncation being
load-bearing.

We then quantify why the natural conditional route does not close.
Writing $\rho(N)=\mathrm{Var}\,C(N)/V(N)$, the coefficient amplifying
the Chowla correlation $S(h)$ grows like $N/(\AAA(N)\log N)$ and is
nonnegative, so a hypothesis $|S(h)|\le\varepsilon$ gives $\rho\to1$
only for $\varepsilon=o(\log N/N)$ --- a factor $N/\log N$ beyond
Chowla --- and no such bound is true, because the measured absolute
budget already exceeds $1$ by an order of magnitude and grows
(Proposition~\ref{prop:W}). The only surviving route is signed
cancellation across $h$, and we state that as an open question.

Finally we delimit what can be measured about $C(N)$ at all. Replacing
$\mu$ by an arbitrary sign pattern on its own support leaves $V(N)$
unchanged, so any estimator reproduced under that replacement is not
measuring $\mu$ (Lemma~\ref{lem:coin}); in particular no centred
estimator calibrates $\rho$. The obstruction is not
information-theoretic --- estimators reading multiplicativity separate
$\mu$ from a coin cheaply, and the field determines $\mu$ by an
invertible additive convolution (Proposition~\ref{prop:coindisc}) ---
but the constructive inverse is numerically dead, its filter growing
like $1.92^{m}$.

Net progress toward Goldbach: zero. What is offered is the exact
facts, and the precise reason two natural routes to $C(N)=o(N)$ are
not routes.
\end{abstract}

% =====================================================================
\section{Introduction}

\subsection{The reduction}

Let $N$ be a large even integer, $\Lambda$ the von Mangoldt function,
$\mu$ the M\"obius function,
$\tilde r(N)=\sum_{n<N}\Lambda(n)\Lambda(N-n)$, and
\begin{equation}\label{eq:constants}
  \AAA(N)=\prod_{q\,\nmid\,N}\Bigl(1-\frac{1}{q(q-1)}\Bigr),
  \qquad
  \SS(N)=2\prod_{p>2}\Bigl(1-\tfrac{1}{(p-1)^2}\Bigr)
          \prod_{\substack{p\mid N\\ p>2}}\Bigl(1+\tfrac{1}{p-2}\Bigr),
\end{equation}
$q$ and $p$ denoting primes. Huang and Li \cite{HL} proved that binary
Goldbach for all sufficiently large even $N$ follows from
$EH_\mu(N^{\theta'})$ for a single $\theta'>1/2$. The companion paper
\cite{P1} evaluates that hypothesis: one half of it is an
unconditional theorem, and the other half is, by an unconditional
identity, equivalent to Huang--Li's equation~(22) --- which yields the
asymptotic $\tilde r(N)\sim\SS(N)N$ exactly when
\begin{equation}\label{eq:C}
  C(N)\;=\;\sum_{n<N}\Lambda(n)\mu(N-n)\;=\;o(N).
\end{equation}
This paper is about \eqref{eq:C}.

The constant $\AAA(N)$ is Huang--Li's equation~(7), and it enters
their argument through the triangle inequality
$|C(N)|\le\sum_{n<N}\Lambda(n)\mu^2(N-n)$, whose right-hand side
counts squarefree shifted primes. Proposition~\ref{prop:V} evaluates
exactly that count. So the deficiency $1-\AAA(N)$ in their lower bound
is one minus the local density of the second moment of $C$, and the
two sides of the problem meet at a constant neither names.

\subsection{Results}

Section~\ref{sec:exact} gives the exact facts:
Proposition~\ref{prop:V} (the scale), Lemma~\ref{lem:MP} (the exact
aggregate identity). Section~\ref{sec:chowla} gives the obstruction to
the conditional route (Proposition~\ref{prop:W}).
Section~\ref{sec:controls} gives the two controls that delimit
measurement (Lemma~\ref{lem:coin}, Proposition~\ref{prop:coindisc}).
Section~\ref{sec:margin} records the size of the margin the problem
demands.

\subsection{What is not claimed}\label{sec:notclaimed}

\begin{itemize}
\item \textbf{No law for $C(N)$ is conjectured here.} In particular no
distributional statement of the form
$C(N)=m(N)+\sqrt{V(N)}\,G(N)$ is asserted.

\item \textbf{No level or rate for $\rho=\mathrm{Var}\,C/V$ is
quoted.} By Lemma~\ref{lem:coin} the centred estimator cannot
distinguish $\mu$ from a sign pattern, so any such figure calibrates
nothing.

\item \textbf{The decoupling is not assumed.} Passing from
$\mathrm{OffDiag}(N)$ to $\sum_{h\ne0}c(h)S(h)$ replaces values
sampled on a sparse set of primes by an average over all integers;
that step is not algebra and is not assumed anywhere below.
Proposition~\ref{prop:W} shows the conditional route fails
\emph{even granting it}.

\item \textbf{No progress toward Goldbach.} Nothing here bounds
$C(N)$.
\end{itemize}

% =====================================================================
\section{The wall, exactly}\label{sec:exact}

\subsection{The exact scale}

\begin{proposition}[the exact scale]\label{prop:V}
Let
\[
  V(N)=\sum_{v<N}\mu^2(v)\,\Lambda(N-v)^2,
  \qquad
  W(N)=\sum_{w<N}\Lambda(w)^2 .
\]
Then $V(N)=W(N)\,\AAA(N)\,(1+o(1))$, and consequently
$V(N)\sim\AAA(N)\,N\log N$.
\end{proposition}

\begin{proof}
$\Lambda(w)^2$ is supported on prime powers with weight $(\log p)^2$,
so with $w=N-v$,
\[
  V(N)=\sum_{p^k<N}(\log p)^2\mu^2(N-p^k)
      =\sum_{p<N}(\log p)^2\mu^2(N-p)+O(\sqrt N\log^2N),
\]
a $(\log p)^2$-weighted count of \emph{squarefree shifted primes}. Its
local density at a prime $q$ is $1$ when $q\mid N$ --- there
$q^2\mid N-p$ forces $q\mid p$, hence $p=q$, a single term --- and
$1-1/\varphi(q^2)=1-1/(q(q-1))$ when $q\nmid N$, since the bad class
$p\equiv N\pmod{q^2}$ is then a unit class. The count is then Mirsky's
theorem on squarefree values of shifted primes \cite{Mir49}, and
partial summation converts it to the stated asymptotic.
\end{proof}

\begin{note}[the local factor is $\AAA$, not $\SS$]\label{note:AvsS}
The singular series $\SS(N)$ is Hardy--Littlewood's and is correct for
the Goldbach \emph{count}; the right object for the second moment of
$C(N)$ is $\AAA(N)$. Both depend on $N$ only through the set of primes
dividing $N$ --- though in opposite ways, $\SS$ through a product
\emph{over} those primes and $\AAA$ through a product over the primes
\emph{not} among them --- which is why the substitution is easy to
make and hard to see. One reconstruction error of exactly this kind
divides by $W(N)$ where the definition calls for $V(N)$; the two
differ by exactly $\AAA(N)$.
\end{note}

\begin{measurement}[the local factor, discriminated]\label{meas:V}
The statistic, stated: for a candidate $c\in\{\AAA,\SS\}$ put
$z_c(N)=\bigl(V(N)/W(N)\bigr)/c(N)$; the figure quoted is
$\mathrm{sd}\bigl(z_c/\overline{z_c}\bigr)$, i.e. the candidate is
rescaled to the measured mean so that only its shape in $N$ is judged.
Over even $N\in[10^5,\,1.6\cdot10^7]$ it is $0.000323$ for $\AAA$
against $0.245235$ for $\SS$, a factor of $759.3$. The reading has to
be stated: subtracting a least-mean multiple rather than dividing
gives $0.000317$ and $0.299040$ instead. The lower cutoff is not
cosmetic --- the figure for $\AAA$ rises as the cutoff falls, reading
$0.000346,\,0.000398,\,0.000470,\,0.000529$ at cutoffs
$5\cdot10^4,\,10^4,\,10^3,\,10^2$, while the figure for $\SS$ is
unaffected, being dominated by $\SS$'s wrong shape rather than by
noise.

Dividing through by $W(N)$ removes the analytic factor exactly, so no
asymptotic for $\sum(\log p)^2$ is assumed: the ratio $V(N)/W(N)$
against $\AAA(N)$ has mean $1.0000002$ with standard deviation
$0.0001659$ in the top octave, and at $N=4\cdot10^6$ the ratio $W/V$
is $1.270800$ against the predicted $1/\AAA(N)=1.270204$. By cells,
with cell $j$ the even $N$ divisible by $2$ and by the first $j-1$ of
$3,5,7,11,13$:
\[
\setlength{\arraycolsep}{3pt}
\begin{array}{r|cccccc}
 \text{modulus} & 2 & 6 & 30 & 210 & 2310 & 30030\\
 \text{count} & 7950000 & 2650000 & 530000 & 75714 & 6883 & 529\\\hline
 \overline{V/(W\AAA)} & 1.000000475 & 1.000000147 & 1.000000010
   & 0.999999570 & 0.999999037 & 0.999999646
\end{array}
\]
(\texttt{lab\_second\_moment.py}).
\end{measurement}

\begin{measurement}[the second-order form]\label{meas:secondorder}
The asymptotic $V\sim\AAA N\log N$ converges slowly and should be
quoted in its second-order form. The measured ratio
$V/(\AAA N\log N)$ is $0.927137,\,0.933409,\,0.939701$ at
$N=10^6,\,4\cdot10^6,\,1.6\cdot10^7$. Partial summation from
$\theta(x)\sim x$ gives $\sum_{p\le x}(\log p)^2\sim x\log x-x$, so
the second-order form is $\AAA(N)\,(N\log N-N)$, a factor
$1-1/\log N$ smaller --- $0.939716$ at the top, against the measured
$0.939701$. Against that form the ratio reads
$0.999482,\,0.999134,\,0.999984$ at the same three $N$
(\texttt{lab\_second\_moment.py}).
\end{measurement}

\subsection{The aggregate second moment, exactly}

\begin{lemma}[the second-moment identity]\label{lem:MP}
Fix $X$ and let
\[
  \widehat C(N)\;=\;\sum_{\substack{n+v=N\\ n\le X,\ v\le X}}
    \Lambda(n)\,\mu(v)
\]
be the truncated convolution, whose support runs to $2X$. Then
\[
  \sum_{N\le2X}\widehat C(N)^2\;=\;\sum_{|h|<X}M(h)\,P(h),
\]
where
\[
  M(h)=\sum_{v,\,v+h\le X}\mu(v)\mu(v+h),
  \qquad
  P(h)=\sum_{w,\,w+h\le X}\Lambda(w)\Lambda(w+h).
\]
\end{lemma}

\begin{proof}
Expanding the square as a double sum over $(n,v)$ and $(n',v')$ with
$n+v=n'+v'$ and summing over \emph{all} $N$ leaves the single
constraint $n-n'=v'-v=:h$, with $(n,n')$ and $(v,v')$ otherwise
ranging independently over $[1,X]^2$. Collecting the $\Lambda$-pairs
and the $\mu$-pairs separately gives the two factors.
\end{proof}

\begin{note}[the truncation is load-bearing]\label{note:trunc}
The truncation on the left is necessary and it is easy to lose. If one
writes instead $\sum_{N\le X}C(N)^2$ with $C$ the untruncated
convolution, the three surviving indices $(n,v,h)$ are constrained by
$n+v\le X$ --- a \emph{simplex} --- while $M(h)P(h)$ ranges them over
the \emph{box} $[1,X]^2$ and so also counts every pair with
$X<n+v\le2X$. The two sides then differ by a factor near $1.57$ at $X$
from $800$ to $3200$, and the ratio does not tend to $1$. The form of
Lemma~\ref{lem:MP} is exact to machine precision at every $X$ tested,
to $X=1.6\cdot10^7$.
\end{note}

The identity is exact and unconditional, and it makes the aggregate
size of the wall a statement about two shifted correlations, one of
which ($P$) is the Hardy--Littlewood prime-pair sum and the other of
which ($M$) is the binary Chowla sum.

% =====================================================================
\section{What the excess is, and why Chowla does not control it}%
\label{sec:chowla}

Write $\rho(N)=\mathrm{Var}\,C(N)/V(N)$, which is $1$ if the $\mu(v)$
on the surviving support are replaced by independent signs. Pointwise,
$C(N)^2=V(N)+\mathrm{OffDiag}(N)$ with
\begin{equation}\label{eq:offdiag}
  \mathrm{OffDiag}(N)
  =\sum_{h\ne0}\ \sum_{\substack{p,\,p+h\ \text{prime}}}
    (\log p)(\log(p+h))\,\mu(N-p)\mu(N-p-h),
\end{equation}
which is algebra. Passing from there to $\sum_{h\ne0}c(h)S(h)$ is
\emph{not} algebra: it replaces the values $\mu(N-p)\mu(N-p-h)$,
sampled on the sparse set of $p$ with $p$ and $p+h$ both prime, by an
average over all $u$. We do not assume that decoupling. What we record
is that even granting it, the conditional route does not close.

Throughout this section
\begin{equation}\label{eq:cdef}
  c(h)\;=\;\sum_{\substack{p'-p=h\\ p,p'\le X}}(\log p)(\log p'),
\end{equation}
the sum being over \emph{primes}, not prime powers. This fixes the
meaning of $c$ once; the alternative
$c(h)=\sum_n\Lambda(n)\Lambda(n+h)$ differs by $1+O(N^{-1/2})$, since
$\psi(x)-\theta(x)\asymp\sqrt x$ and the relevant quantity is squared,
and nothing below turns on the choice.

\begin{proposition}[the amplification]\label{prop:W}
Let $\Gamma(N)=\bigl(\sum_{h\ne0}c(h)\bigr)/V(N)$ with $c$ as in
\eqref{eq:cdef} and $X=N$. Then
\[
  \Gamma(N)\;\sim\;\frac{N}{\AAA(N)\log N},
\]
and consequently a hypothesis $|S(h)|\le\varepsilon$ yields nothing
better than $|\rho-1|\le\varepsilon\,\Gamma(N)$. Because $c(h)\ge0$
that inequality is sharp, so such a hypothesis can give $\rho\to1$
only for $\varepsilon=o(1/\Gamma(N))=o(\log N/N)$.
\end{proposition}

\begin{proof}
The numerator needs no correlation input: summing $c(h)$ over all
$h\ne0$ is exactly $\theta(N)^2-\sum_{p<N}(\log p)^2$, which is
$N^2(1+o(1))$ by the prime number theorem. The denominator is
$\AAA(N)N\log N(1+o(1))$ by Proposition~\ref{prop:V}. Dividing gives
the asymptotic. The inequality $|\rho-1|\le\varepsilon\Gamma$ is the
triangle inequality applied to $\sum_{h\ne0}c(h)S(h)$ after division
by $V$, and it is sharp because every $c(h)$ is nonnegative, so no
cancellation among the coefficients is available to a hypothesis that
controls only $|S(h)|$.
\end{proof}

Two consequences. First, Chowla's conjecture asserts $S(h)=o(1)$ for
fixed $h$ --- and the strongest unconditional results in that
direction, the logarithmically averaged two-point theorem of
\cite{Tao16} and the shift-averaged bound of \cite{MRT15}, give
smallness of exactly that order. The strength the display needs is
$S(h)=o(\log N/N)$, so the gap is a factor $N/\log N$. Second --- and this is what closes the
absolute-value route rather than merely costing it a power of $\log$
--- the true $S$ overspends that budget already. Define the
\emph{absolute budget}, which is exactly what the triangle inequality
spends,
\begin{equation}\label{eq:budget}
  \mathcal B(X)=\frac{1}{V(X)}\sum_{0<|h|<X}c(h)\,\bigl|S(h)\bigr|,
  \qquad
  S(h)=\frac{1}{X-|h|}\!\!\sum_{|h|<u\le X}\!\!\mu(u)\,\mu(u-|h|).
\end{equation}

\begin{measurement}[the budget is already overspent]\label{meas:budget}
Measured, $\mathcal B(X)=13.3,\,17.3,\,23.1,\,30.5$ at
$X=2\cdot10^4,\ 4\cdot10^4,\ 8\cdot10^4,\ 1.6\cdot10^5$: above $1$ by
an order of magnitude, and growing. Normalising $S$ by $X$ rather than
by the number $X-|h|$ of terms gives $7.9,\,10.3,\,13.8,\,18.2$ ---
smaller, same conclusion. Both are reported because the phrase
``binary Chowla sum'' does not fix the normalisation, and a budget is
not interpretable until it does. The amplification itself measures
$\Gamma=1.5128\cdot10^{3},\ 1.8412\cdot10^{4},\ 3.5759\cdot10^{5}$ at
$N=10^4,\,1.6\cdot10^5,\,4\cdot10^6$, with
$\Gamma\log N/N=1.3933,\,1.3789,\,1.3590$ against the predicted
$1/\AAA(N)=1.270204$
(\texttt{audit\_amplification.py}).
\end{measurement}

Hence no bound on $|S(h)|$ alone can give $\rho\to1$. The only route
is \emph{signed cancellation across $h$}, which is a different object
from Chowla's smallness and from the averaged absolute bound of
\cite{MRT15}. Naming what supplies that cancellation is open, and we
state it as such.

% =====================================================================
\section{Two controls, and what they delimit}\label{sec:controls}

\begin{lemma}[the coin control]\label{lem:coin}
Let $\varepsilon(v)=\pm1$ be arbitrary signs on
$\{v:\mu(v)\neq0\}$ and zero elsewhere. Then $\varepsilon^2=\mu^2$
pointwise, so $V(N)$ is unchanged, and the field
$C_\varepsilon(N)=\sum_v\varepsilon(v)\Lambda(N-v)$ has the same exact
second moment as $C(N)$ for every $N$. Consequently any estimator
whose output is reproduced when $\mu$ is replaced by $\varepsilon$ is
not measuring $\mu$.
\end{lemma}

\begin{proof}
Immediate from $\varepsilon^2=\mu^2$ and the definition of $V$.
\end{proof}

\begin{note}[what the lemma licenses, and what it does not]%
\label{note:coinscope}
The implication runs one way: \emph{reproduced under $\varepsilon$}
$\Rightarrow$ \emph{not measuring $\mu$}. The converse does not follow
and is not proved anywhere here, so a statistic that a sign ensemble
fails to reach has not thereby been shown to measure $\mu$; it has
been shown not to be killed by this particular control. The
distinction is quantitatively large. An i.i.d. sign ensemble is
\emph{narrower} than a random multiplicative one wherever the object
is multiplicative --- resolved to $512$ draws on one statistic, the
i.i.d. spread was $0.6455$ of the multiplicative
(	exttt{audit\_cn\_coin\_deep.py}) --- so significance
read against an i.i.d. ensemble is larger than it should be. Because
narrowness is conservative for killing, every claim this lemma kills
stays killed; it is the other direction that fails. An estimator that
escapes an i.i.d. ensemble has cleared the weakest available control,
and until it has been put against an ensemble sharing $\mu$'s
multiplicativity, the escape is a fact about the control.
\end{note}

Lemma~\ref{lem:coin} invalidates the natural measurement of $\rho$:
the centred estimator returns the same value on $\varepsilon$ as on
$\mu$, so a quoted level for $\rho$ calibrates nothing. \emph{Any
claim about $\mathrm{Var}\,C$ relative to $V$ must first exhibit an
estimator that distinguishes $\mu$ from a sign pattern, and show that
it does.}

\begin{proposition}[the coin obstruction is arithmetic, not
informational]\label{prop:coindisc}
\begin{enumerate}
\item[(i)] Estimators separating $\mu$ from a sign pattern exist and
are cheap, provided they read \emph{multiplicativity} rather than
variability. Since $\mu(2v)=-\mu(v)$ for odd $v$ is an exact identity
while $\varepsilon(2v)$ and $\varepsilon(v)$ are independent signs,
\[
  T(x)\;=\;\frac1x\sum_{v\le x}\mu(v)\mu(2v)
  \;\longrightarrow\;-\frac{4}{\pi^2},
\]
minus the density of odd squarefree integers, while its analogue for
independent signs is $O(x^{-1/2})$.
\item[(ii)] The field determines $\mu$. Writing $M=N-2$ and
$\widehat\Lambda(m)=\Lambda(m+2)$, one has
$C(N)=(\widehat\Lambda*\mu)(M)$ as an additive convolution with
$\widehat\Lambda(0)=\Lambda(2)=\log2\neq0$, so $\widehat\Lambda$ is
invertible and $\mu(M)=\sum_{m<M}a(m)\,C(M-m+2)$ for the inverse
filter $a$.
\end{enumerate}
Hence Lemma~\ref{lem:coin} is not an information-theoretic
obstruction: two different sign patterns give two different fields.
\end{proposition}

\begin{proof}
(i) On odd squarefree $v$, $\mu(2v)=\mu(2)\mu(v)=-\mu(v)$, so
$\mu(v)\mu(2v)=-\mu^2(v)$ there; on even $v$, $\mu(2v)=0$. Summing
gives $-\#\{v\le x:v\ \text{odd squarefree}\}=-(4/\pi^2)x(1+o(1))$.
For independent signs the summands are independent mean-zero, whence
the $O(x^{-1/2})$. (ii) An additive convolution by a sequence with
nonzero value at $0$ is invertible over $\mathbb Q$ by forward
substitution, which defines $a$ uniquely and gives the stated
inversion.
\end{proof}

\begin{measurement}[the discriminator, and the price of inverting]%
\label{meas:disc}
At $x=10^4,\,10^5,\,10^6,\,2\cdot10^6$, against
$-4/\pi^2=-0.405285$,
\[
  T(x)\;=\;-0.405600,\ -0.405270,\ -0.405286,\ -0.405295,
\]
while the worst of twenty independent-sign draws on the same support
is
\[
  0.020600,\ 0.004570,\ 0.001214,\ 0.000935 .
\]
At the top the separation is about $430$ standard deviations of the
sign ensemble.

The constructive route of (ii) is killed by amplification, not by
information. The inverse filter grows geometrically: $|a(m)|$ runs
$1.443,\,2.287,\,2.182,\,20.91,\,512.3,\,3.443\cdot10^{5},\,
1.073\cdot10^{14},\,1.535\cdot10^{28},\,3.146\cdot10^{56}$ at
$m=0,1,2,5,10,20,50,100,200$, a growth of about $1.92^{m}$. Any
numerical use of the inversion is therefore dead at once
(\texttt{lab\_coin\_discriminator.py}).
\end{measurement}

% =====================================================================
\section{The margin the problem demands}\label{sec:margin}

The requirement $C(N)=o(N)$ constrains every $N$, so the figure to
quote is the one at the extreme rather than at a typical $N$. With
$\max|C|\approx a_n\sqrt{V(N)}$ and $a_n$ the Gumbel location for the
number of even $N$ below the point,
\[
  \frac{N}{\max_{N\le X}|C(N)|}\;\approx\;
  \frac{\sqrt N}{a_n\sqrt{\AAA\log N}},
\]
which is $10^{4.466}$ at $N=10^{12}$ and $10^{22.842}$ at $N=10^{50}$
with $a_n=\sqrt{2\log(N/2)}$ and $\AAA=0.787275$, the value of
\eqref{eq:constants} at a generic even $N$
(	exttt{audit\_margin.py}, 	exttt{audit\_constants.py}). The quantity is
therefore expected to be small by a wide margin and to be provably
small by none: the gap between what is expected and what is reachable
is the whole problem.

% =====================================================================
\section{Open questions}

\begin{enumerate}
\item $C(N)=o(N)$ itself.
\item Name the mechanism supplying \emph{signed} cancellation across
$h$ in \eqref{eq:offdiag}. Proposition~\ref{prop:W} shows that no
bound on $|S(h)|$ can substitute for it, and nothing in the literature
supplies it.
\item Exhibit an estimator of $\rho$ that provably distinguishes $\mu$
from a sign pattern in the sense of Lemma~\ref{lem:coin}, and
calibrate it against a random multiplicative ensemble rather than an
i.i.d. one (Note~\ref{note:coinscope}).
\end{enumerate}

% =====================================================================
\section{Summary}

\begin{itemize}
\item Proposition~\ref{prop:V}: $V(N)\sim\AAA(N)N\log N$, with the
local factor $\AAA$ and not the singular series $\SS$; the same
constant is the deficiency in Huang--Li's lower bound.

\item Lemma~\ref{lem:MP}: the aggregate second moment of the truncated
convolution is exactly $\sum_h M(h)P(h)$, unconditionally. The
truncation is load-bearing.

\item Proposition~\ref{prop:W}: the coefficient amplifying $S(h)$
grows like $N/(\AAA\log N)$ and is nonnegative, and the measured
absolute budget already exceeds $1$ by an order of magnitude and
grows. So no bound on $|S(h)|$ gives $\mathrm{Var}\,C\sim V$; only
signed cancellation across $h$ would.

\item Lemma~\ref{lem:coin} and Proposition~\ref{prop:coindisc}: no
centred estimator calibrates $\rho$; discriminating estimators exist
and must read multiplicativity; the constructive inversion is
numerically dead.

\item Net progress toward Goldbach: zero.
\end{itemize}

% =====================================================================
\begin{thebibliography}{99}

\bibitem{HL} J.-J.~Huang and H.~Li, \emph{On the connection between
the Goldbach conjecture and the Elliott--Halberstam conjecture},
arXiv:2005.03811v2 [math.NT], 2022.

\bibitem{Mir49} L.~Mirsky, \emph{The number of representations of an
integer as the sum of a prime and a $k$-th power free integer},
Amer. Math. Monthly \textbf{56} (1949), 17--19.

\bibitem{MRT15} K.~Matom\"aki, M.~Radziwi{\l}{\l} and T.~Tao,
\emph{An averaged form of Chowla's conjecture}, Algebra Number Theory
\textbf{9} (2015), 2167--2196.

\bibitem{P1} \emph{An unconditional bound for a M\"obius-weighted
correlation sum in fixed residue classes, with two consequences for
the Huang--Li reduction of Goldbach to Elliott--Halberstam},
companion paper.

\bibitem{Tao16} T.~Tao, \emph{The logarithmically averaged Chowla and
Elliott conjectures for two-point correlations}, Forum Math. Pi
\textbf{4} (2016), e8.

\end{thebibliography}

% =====================================================================
\appendix
\section{Reproducibility}

\begin{center}
\begin{tabular}{lll}
\hline
Script & Result file & Supports\\
\hline
\texttt{lab\_second\_moment.py} & \texttt{lab\_second\_moment.txt}
 & Measurements~\ref{meas:V},~\ref{meas:secondorder}\\
\texttt{audit\_amplification.py} & \texttt{audit\_amplification.txt}
 & Measurement~\ref{meas:budget}\\
\texttt{lab\_coin\_discriminator.py} &
 \texttt{lab\_coin\_discriminator.txt} & Measurement~\ref{meas:disc}\\
\hline
\end{tabular}
\end{center}

Every measurement above states the range on which it was taken, and
the statistic and field it was taken of. No computation is used as a
premise of Proposition~\ref{prop:V}, Lemma~\ref{lem:MP},
Proposition~\ref{prop:W} or Proposition~\ref{prop:coindisc}, each of
which is proved.

\end{document}
